% 付録：階差型・総和型から微積分と近似計算へ

\section{和と差から微積分へ}\label{節：和と差から微積分へ}

\sectionref{節：階差型}は差から元の\bookterm{sequence}{数列}を求め,\sectionref{節：総和型}は和から差を取り出した.
この逆向きの二つの操作を,微分と積分につなげよう.
隣の項との差は,一歩進んだときの変化である.
変化を取り出す操作と,その変化を足して元へ戻す操作を,同じ流れの中で見ていく.

\subsection{和と差の対応}
以下では,最初の項を$a_0$と書くことがある.
例えば平方の列$0,1,4,9,\ldots$から隣り合う差を取ると,$1,3,5,\ldots$になる.
\bookterm{sequence}{数列}$a_0,a_1,a_2,\ldots$の隣り合う項の差
\[
  \Delta a_n=a_{n+1}-a_n
\]
を一\bookterm{difference}{階差}分という.
本編の階差に$\Delta$という記号と「\bookterm{difference-operator}{差分}」の名前を付けた.
たとえば$a_n=n^2$なら
\[
  \Delta a_n=2n+1,\qquad \Delta^2a_n=2.
\]
二次式,一次式,定数の順になる.$\Delta^2a_n=\Delta a_{n+1}-\Delta a_n$である.

\sectionref*{節：特性方程式型}の別解でも,添字を一つずらして引くことで定数項$q$を消し,
\[
  \Delta a_{n+1}=p\,\Delta a_n
\]
という等比型を作った.
これも\bookterm{difference-operator}{差分}で扱いやすい量を取り出す例である.

逆に,\bookterm{difference-operator}{差分}を足し合わせれば元の\bookterm{sequence}{数列}を求め直せる.
\[
  \begin{aligned}
    a_n-a_0
    &=(a_1-a_0)+\cdots+(a_n-a_{n-1})\\
    &=\sum_{k=0}^{n-1}\Delta a_k.
  \end{aligned}
\]
途中の項が消える.本編の\bookterm{difference}{階差}型の和は離散的な積分,総和型の$S_{n+1}-S_n=a_{n+1}$はその逆になる.図で微積分との対応を見る.
\BookFigDifference

この見方は,\sectionref{節：階差等比複合型}で引く多項式$g(n)$を探したことにもつながる.
特に$p=1$なら,$g(n+1)-g(n)=f(n)$を満たす式を探す問題になる.
たとえば$S_n=\sum_{k=1}^n k^2$については
\[
  S_n-S_{n-1}=n^2,\qquad S_0=0.
\]
二次式の差を逆向きに戻すので,三次式
$S_n=An^3+Bn^2+Cn+D$を候補にする.
その差は
\[
  S_n-S_{n-1}
  =3An^2+(-3A+2B)n+(A-B+C).
\]
係数比較から$\displaystyle A=\frac{1}{3}$, $\displaystyle B=\frac{1}{2}$, $\displaystyle C=\frac{1}{6}$となり,$S_0=0$から$D=0$である.
同じ初期値と差を満たすので
\[
  S_n=\frac{n(n+1)(2n+1)}6
\]
が確かめられる.
和の公式は,差を取ると目的の式になるものを探して作れる.

\subsection{二項係数で多項式を表す}\label{節：二項係数で多項式を表す}
平方和を,今度は\sectionref*{節：二項定理}の二項係数で求める.
ここでは$n,k,N$は非負整数とし,
\[
  \binom nk=\frac{n!}{k!(n-k)!}\quad(0\le k\le n)
\]
とする.$0!=1$, $k>n$では$\binom nk=0$と約束する.

\bookstep{例題：二項係数で平方和を求める}
$n^2=n(n-1)+n$だから,
\[
  n^2=2\binom n2+\binom n1
\]
と表せる.
部品の差は
\[
  \binom{n+1}{3}-\binom n3=\binom n2\quad(n\ge0)
\]
となる.$n\ge3$では階乗表示を展開して確かめられ,$n=0,1,2$も上の規約で成り立つ.
同じように,$\binom{n+1}{2}-\binom n2=\binom n1$である.
$j=0$から$N-1$まで足すと途中が消え,従って
\BookMathLayout{\[
    \sum_{j=0}^{N-1}j^2=2\binom N3+\binom N2
    =\frac{N(N-1)(2N-1)}6.
  \]}{\[
    \begin{aligned}
      \sum_{j=0}^{N-1}j^2&=2\binom N3+\binom N2\\
      &=\frac{N(N-1)(2N-1)}6.
    \end{aligned}
  \]}
$N=n+1$とすれば,先ほどの$1^2+\cdots+n^2$の公式に戻る.
二項係数なら,添字を一つ下げるだけで差を表せる.

\bookstep{一般の添字で確かめる}
$n+1$個から$k$個を選ぶとき,特定の一個を選ぶかどうかで分ければ,
\[
  \binom{n+1}{k}=\binom nk+\binom n{k-1}\quad(k\ge1)
\]
である.従って
\[
  \Delta\binom nk=\binom n{k-1}\quad(k\ge1)
\]
となる.$\Delta$は前進の差で,$\binom n0=1$の差は別に零となる.

逆に差を足すと,途中の項が消えるので,
\[
  \sum_{j=0}^{N-1}\binom jk=\binom N{k+1}
  \quad(N,k\ge0)
\]
となる.$N=0$では空和$0$,$k=0$では和が$N$となり,右辺と一致する.
\BookFigBinomialDifference

\begin{bookpoint}
  \textbf{式の部品を選ぶ}\par
  目的の操作が簡単になる\bookterm{basis}{基底}を選ぶ.
  \sectionref*{節：隣接3項型}の見方が,多項式の表示にも現れた.
\end{bookpoint}
実際,$\binom n0,\ldots,\binom nd$はそれぞれ次数が$0,\ldots,d$で,
最高次係数が非零である.高い次数から係数を合わせれば,
$d$次以下の多項式を一意にその\bookterm{combination}{線形結合}で表せる.
三乗和は\bookproblemref{practice-connections}{7}{接続演習問7}で試そう.


\subsection{二方向の階差}\label{節：二方向の階差}
非負整数$n,m$で数表の値$a_{n,m}$を表す.
\sectionref*{節：階差型}の和と差を,横と縦の二方向で使おう.

\bookstep{例：距離の二乗と階差}
原点から点$(n,m)$までの距離は
\[ a_{n,m}=\sqrt{n^2+m^2} \]
である.
根号を取り除くため$b_{n,m}=a_{n,m}^2$とすると,
一方向の差は$b_{n+1,m}-b_{n,m}=2n+1$となり,$m$によらない.
従って,さらにもう一方向へ差を取れば零になる.

\bookstep{零の差から一定の差へ}
この計算を広げ,境界の値
\[
  b_{n,0}=n^2,\qquad b_{0,m}=m^2
\]
と,一定の実数$c$について
\BookMathLayout{\[
    b_{n+1,m+1}-b_{n+1,m}-b_{n,m+1}+b_{n,m}=c
  \]}{\[
    \begin{aligned}
      b_{n+1,m+1}-b_{n+1,m}\\
      {}-b_{n,m+1}+b_{n,m}=c
    \end{aligned}
  \]}
とする.$b_{n+1,m+1}$は添字の和が小さい三つの値から決まるので,$n+m$についての帰納法で境界から一意に求まる.

横方向の差$d_{n,m}=b_{n+1,m}-b_{n,m}$を追うと,
\[
  d_{n,m+1}-d_{n,m}=c,\qquad d_{n,0}=2n+1
\]
なので$d_{n,m}=2n+1+cm$である.これを横に足し戻せば,
\[
  b_{n,m}=n^2+m^2+cnm
\]
となる.$c$が一定なので,\textbf{二方向の\bookterm{difference}{階差}を逆順に足し戻せた}.

\bookstep{数表と図形}
二つの単位ベクトル$\boldsymbol u,\boldsymbol v$のなす角を$\theta$とすると,
$c=2\cos\theta$の場合には
\[
  b_{n,m}=|n\boldsymbol u+m\boldsymbol v|^2
\]
と読める.数表なら$c$は任意の実数でよいが,この距離の二乗という解釈には$-2\le c\le2$が必要になる.

\bookstep{例：格子経路を数える}
一方,格子点$(0,0)$から$(n,m)$へ,右と上へ一歩ずつ進む道の数を$A_{n,m}$とする.
最後の一歩を分ければ,
\[
  A_{n,m}=A_{n-1,m}+A_{n,m-1}\quad(n,m\ge1)
\]
であり,境界では$A_{n,0}=A_{0,m}=1$である.
\BookFigLatticeRecurrence
下の数表は,上と左の値を足して作ったものである.
\[
  \begin{array}{c|rrrr}
    m\backslash n&0&1&2&3\\\hline
    0&1&1&1&1\\
    1&1&2&3&4\\
    2&1&3&6&10\\
    3&1&4&10&20
  \end{array}
\]
全$n+m$歩から上へ進む$m$歩の位置を選ぶので,
\[
  A_{n,m}=\binom{n+m}{m}
\]
である.最後の一歩の分類は,二項係数の加法公式でもある.
境界と帰納法による確認は\bookproblemref{practice-connections}{3}{接続演習問3}へ.

\bookstep{二つの例の共通点}
このように二つの添字をもつものを\textbf{\bookterm{double}{二重数列}}という.
一変数で使った階差と帰納法を二方向へ広げた見方である.

\begin{bookpoint}
  \textbf{二つの添字と初期条件}\par
  \bookterm{initial}{初期条件}となる境界と更新順序を確かめる.
  二方向の規則を別々に与えるなら,同じ位置を二通りに計算して一致する必要がある.
\end{bookpoint}
格子経路と二項係数の資料は参考文献のNIST DLMF §26.3へ.


\subsection{差分商と微分}
間隔$h>0$で平方を記録すると,$y(t)=t^2$について
\[
  \frac{y(t+h)-y(t)}h
  =\frac{(t+h)^2-t^2}h=2t+h.
\]
となり,$h\to0$で傾き$2t=y'(t)$に近づく.一般に$t_n=nh$での値を$a_n$とすると,単位時間あたりの変化は
\[
  \frac{a_{n+1}-a_n}{h}
\]
である.
微分可能な関数$y$について,\textbf{時刻$t$を固定}して$h$を小さくすると,\footnote{基点も動かす近似では,例えば$y'$がある閉区間で連続なら,その区間に収まる$t,t+h$について差分商と$y'(t)$の差が一様に小さくなる.ここで使う$y(t)=t^2$はこの条件を満たす.}
\[
  \frac{y(t+h)-y(t)}h\longrightarrow y'(t)
\]
となる.格子点$t_n$での傾きもこの比で近似できるが,基点も動かす場合の一様な近似には追加条件が要る.

逆に,微分方程式$y'=F(t,y)$を数値で解くには,今の点での傾きから
\[
  a_{n+1}=a_n+hF(t_n,a_n)
\]
と次の値を予測できる.
これをオイラー法という.接線の方向へ一歩ずつ進む\bookterm{recurrence}{漸化式}で曲線を近似する.

例えば$y'=-y$, $y(0)=1$の解は$e^{-t}$で,$0$に近づく.
オイラー法では
\[
  a_{n+1}=(1-h)a_n,\qquad a_0=1
\]
だから,$n\ge1$では$a_n=(1-h)^n$である.最初の値は$a_0=1$のまま保つ.
本編の等比型で調べたように,$0$への収束条件は$|1-h|<1$, すなわち$0<h<2$になる.
$0<h<1$では正のまま減少し,$h=1$なら一回で$0$になる.
$1<h<2$では符号を交互に変えながら近づくが,$h=2$では$1,-1,1,-1,\ldots$と繰り返して収束しない.
さらに$h>2$では絶対値が増大する.
本当の解が減衰しても,刻み幅次第で近似の列は増大する.等比型の公比は,数値計算の安定性も教える.

また,正の時刻$t$を$N$等分して$\displaystyle h=\frac{t}{N}$とすれば,時刻$t$における近似値は
\[
  a_N=\left(1-\frac{t}{N}\right)^N
  \longrightarrow e^{-t}.
\]
この例では,刻みを細かくした近似の収束を指数関数の\bookterm{limit}{極限}で確かめられる.一般の保証には関数や刻み幅の条件が要る.

\subsection{ニュートン法の誤差}
\sectionref{節：ニュートン法}では,\BookMathLayout{}{\\}%
$a_{n+1}=\frac12\left(a_n+\dfrac{2}{a_n}\right)$の$\sqrt2$への誤差がほぼ二乗された.
一般の\bookterm{newton}{ニュートン法}でも,誤差が二乗の定数倍以下になる条件を調べよう.

方程式$g(x)=0$に対する\bookterm{newton}{ニュートン法}の更新は
\[
  x_{n+1}=x_n-\frac{g(x_n)}{g'(x_n)}
\]
である.
$g(\alpha)=0$, $g'(\alpha)\ne0$とし,根$\alpha$の近くで$g$は二回連続微分可能とする.
本編のずれの置き換えと同じく,$e_n=x_n-\alpha$とおくと
\[
  e_{n+1}
  =\frac{g'(x_n)e_n-g(x_n)}{g'(x_n)}.
\]
分子は,接線による変化と実際の変化との差である.
二階微分$g''$の絶対値がこの近くで$M$以下なら
\[
  \begin{aligned}
    g'(x_n)e_n-g(x_n)
    &=\int_\alpha^{x_n}\{g'(x_n)-g'(t)\}\,dt,\\
    |g'(x_n)e_n-g(x_n)|
    &\le \frac{M}{2}|e_n|^2.
  \end{aligned}
\]
平均値の定理で$|g'(x_n)-g'(t)|\le M|x_n-t|$と評価して積分した.$x_n<\alpha$でも絶対値で同じ評価になる.
また,$g'(\alpha)\ne0$と連続性により,十分近くでは$|g'(x_n)|\ge m$となる正の定数$m$を選べる.
したがって
\[
  |e_{n+1}|\le\frac{M}{2m}|e_n|^2.
\]
十分小さな誤差は二乗の定数倍以下に縮むので,近くにとどまり急速に根へ収束する.
本編の$\sqrt2$の例で得た正確な式
\[
  e_{n+1}=\frac{e_n^2}{2a_n}
\]
は,その具体例である.

前節の\bookterm{fixed}{不動点}の見方でも確かめられる.
更新関数$\displaystyle N(x)=x-\frac{g(x)}{g'(x)}$について
\[
  N'(\alpha)
  =\left.\frac{g(x)g''(x)}{g'(x)^2}\right|_{x=\alpha}=0.
\]
つまり,根のところでは一次のずれが消え,次の誤差は二乗の定数倍以下になる.
二次の項まで消えると,さらに速く近づく場合がある.\footnote{例えば$g(x)=x+x^3$の根$0$では,$N(x)=\dfrac{2x^3}{1+3x^2}$となり,誤差は三次の大きさになる.$g(x)=x$なら一回で根に移る.}
等比型の一定倍率と違い,\bookterm{newton}{ニュートン法}では誤差が小さいほど残る割合も小さくなる.

接線の傾きを毎回微分で求めるのが難しいなら,過去二点を結ぶ割線の傾きで代用できる.
その更新は
\[
  x_{n+1}
  =x_n-g(x_n)\frac{x_n-x_{n-1}}{g(x_n)-g(x_{n-1})}
\]
となり,割線法と呼ばれる.
分母が零でない範囲で用いる.
直前の二項を状態にする点は前節と共通する.\bookterm{recurrence}{漸化式}は,近似の規則を作り収束を調べる道具にもなる.
